%% Discussion (plan steps 10.6 and 7.4). Every number is a macro from tables/numbers.tex.
\paragraph{What a defined comparison of releases looks like.} One implication follows from RQ1 alone and does not depend on the detection results. Because flow-level recall weights each conversation by the number of records its exporter wrote (Eq.~\eqref{eq:flowrecall}), and because that number varies between releases by class (Table~\ref{tab:rq1-fpc}), a flow-level comparison between two releases of \cicids{} compares different estimands for every attack class except PortScan, whatever the detectors do. RQ5 shows that the difference is large enough to change published conclusions. A comparison of two releases is defined when it holds four things fixed: the conversations scored, the labels they carry, the split that sends them to training or test, and the unit at which they are counted. When all four are fixed, the gap that remains belongs to the exporter's features and segmentation as they reach the detector, which is what a claim that one release is better should be about. Fixing the labels includes a choice that is easy to miss: how a record without attack payload is labelled. Under the schedule label it is an attack, under a payload rule it is benign, and on \cicids{} that choice alone decides whether the RQ3 drop and the RQ4 residual appear (Sections~\ref{sec:rq3} and \ref{sec:rq4}). The conversation unit is not neutral, though. Under the ``any'' rule it turns an end-of-connection record that the schedule labels as attack into a false alarm on benign traffic (Section~\ref{sec:rq4}), and a byte-weighted rule removes that effect. A per-conversation score therefore needs its aggregation rule stated, as a per-flow score needs its split.

\paragraph{A reporting checklist.} Table~\ref{tab:checklist} turns these results into items that a producer of a flow release or an evaluator can apply without our code. Each item names the result that makes it necessary.

\begin{table*}[t]
\caption{Reporting checklist for flow releases and for evaluations that use them. Evidence: the result in this paper that makes the item necessary.}\label{tab:checklist}
\small
\begin{tabularx}{\textwidth}{@{}l>{\raggedright\arraybackslash}p{5.6cm}L@{}}
\toprule
Item & What to report & Evidence \\
\midrule
\multicolumn{3}{@{}l}{\emph{For producers of a flow release}} \\
P1 & The exporter, its version, every timeout and expiry setting (active, idle, FIN and RST expiry), and any filter applied after export. & Two releases of the same packets that differ in one expiry rule and the filter it triggers differ in records per conversation for most attack classes and in flow-level F1 (Section~\ref{sec:rq3}). \\
P2 & The five-tuple and the start and end timestamps of every record, at their native resolution. & Without them records cannot be grouped into conversations; the minute clock of the original \cicids{} release already forces a mapping step (Section~\ref{sec:key}). \\
P3 & The labelling rule (schedule, windows, what is dropped, how records without attack payload are labelled), and a flag wherever a record's label departs from it. & Scored against their own labels, the seven \cicids{} releases agree on the model ranking with $W$ of \RQfourOwnScoreWFlow{}, against \RQfourWFlow{} with one schedule truth (Section~\ref{sec:rq4}); under payload-aware labels the RQ3 drop disappears (Section~\ref{sec:rq3}). \\
\midrule
\multicolumn{3}{@{}l}{\emph{For evaluators}} \\
E1 & Split by conversation (endpoint pair, protocol and a gap $T$), never by record, and state whether larger units such as attack campaigns straddle the split. & The leakage of a random flow split differs by a factor of more than \RQoneLeakFactor{} between releases (Section~\ref{sec:rq1}); the conclusions do not change for $T$ from \NsweepCicStableMin{} to \setnum{120} seconds, and every \cicids{} attack episode straddles the split (Section~\ref{sec:rq4}). \\
E2 & Score per conversation as well as per record, on the conversations both releases contain, and state the aggregation rule. & Per conversation, \RQfiveReverse{} published contrasts reverse and \RQfiveContradicted{} claims of no difference fail (Section~\ref{sec:rq5}); the ``any'' rule creates the residual that a byte-weighted rule removes (Section~\ref{sec:rq4}). \\
E3 & Per-class F1 with the number of independent test units per class and seed: conversations, and episodes where the schedule defines them. & A flow count overstates the number of independent units by the records-per-conversation factor of the release (Table~\ref{tab:rq1-fpc}). \\
E4 & Flag, and leave out of a headline average, any class with fewer than \SetMinTest{} independent units on either side of the split or whose units all straddle it. & Heartbleed and Infiltration, below; \RQfiveNotEstimable{} published contrasts are not estimable for this reason (Section~\ref{sec:rq5}). \\
\bottomrule
\end{tabularx}
\end{table*}

Four classes show how the items work. \emph{Heartbleed} fails items E3 and E4. The whole class is one connection, written as \RQoneHeartbleedRecords{} records by every release except BCCC. Under the conversation-grouped split it is in training in all five seeds and never in test, so no conversation-level F1 exists for it; under a random flow split its records are divided between training and test, and an F1 reported for it measures how well a model recognises sibling records of a connection it was trained on. It should be reported as one unit that cannot be evaluated. \emph{Infiltration} fails items P3, E3 and E4. The schedule we use excludes it, and in the releases that label it all its conversations connect one pair of hosts, so it is at most one independent unit however the records are split; two releases do not contain it. \emph{UNSW-NB15 Backdoor} passes: it has \RQfourBackdoorTest{} test and \RQfourBackdoorTrain{} training conversations in every seed, and \RQfourBackdoorEpTest{} test episodes, of which \RQfourBackdoorEpStrad{} also have conversations in training. Reported with these counts, its F1 can be compared across releases. \emph{CICIDS2017 SSH-Patator} is a borderline case. It passes at the conversation level, with \RQfourSshTestConv{} test conversations per seed, but it is one campaign against one victim, and that campaign straddles every split, so its test conversations are independent of training only at the level of the conversation.

\paragraph{Near-perfect scores.} Near-perfect flow-level scores on \cicids{} are common \cite{maseer2021}. Our results do not show that such scores are inflated by leakage: in RQ3 the release that leaks most gains almost nothing from it, because the sibling record in training shares almost nothing with the record in test. They show something different. A flow-level score on one release refers to a population of records that another release of the same traffic does not have, with each conversation counted as many times as the exporter cut it. Two near-perfect scores on two releases are therefore not two measurements of the same detector quality, and their difference is not evidence that one release is better.

\paragraph{What we did not test.} Janati Idrissi et al.\ report a spread of up to \litnum{8.77} percent in macro-F1 between timeout settings \cite{janatiidrissi2025}. Their design changes the scored flows with every setting (Section~\ref{sec:related}), so the spread can come from the detectors or from the population, and we did not test how much of it the population change alone could produce; that needs a re-extraction of the captures over their grid, which is outside our scope of published releases. The UNSW-NB15 residual of RQ4 is not explained, and a feature ablation would be the next test. Own-label and payload-aware retraining were run on \cicids{} only.
